\section{What an accepted reading permits}
\label{sec:source-admission}

Apple's inventory policy illustrates the distance between reading a sentence
and reconstructing a business. Its reported use of FIFO supports a proposed
cost-formula reading for the inventory and reporting contexts covered by that
policy. It does not reveal the quantities and ages of the actual lots, establish
that no write-down occurred, or supply the authority for every US accounting
treatment. Those are separate inputs to a cost and impairment calculation.
Similarly, a straight-line depreciation policy leaves useful lives, residual
values, components and available-for-use dates to be established.

The source preparation retains six located Apple policy passages: consolidation
of Apple and wholly owned subsidiaries; comparative reclassification; FIFO;
straight-line depreciation; the interim statement's estimates and annual-report
reference; and its 52/53-week fiscal calendar. Each proposed reading carries
the original document hash, extraction hash and exact character slice. These
are company assertions. The consolidation sentence, for example, cannot exclude
every partly owned entity without a control analysis. The calendar sentence
does not make every future fiscal year 52 weeks or make a contractual interest
year 365 days. The exact quotations and proposed predicates are retained in
\path{../methods/source-translations-v1.json}; the meaning explained here is
not an invented human acceptance of them.

A reading progresses through separate decisions. First, extraction review asks
whether the located words and numbers were transcribed faithfully. Translation
review asks whether the proposed predicate preserves their scope, quantifier,
exceptions and reporting context. A model then admits that predicate, or an
explicitly labelled assumption, for a particular purpose and information view.
Accounting eligibility subsequently asks whether a method applies to the actual
event. A faithful transcription of an erroneous company assertion is still
faithful transcription. The four decisions therefore cannot share a single
``approved'' flag.

An explicit presumption in a standard also differs from filling an unknown with
a convenient answer. Article 6 of the retained official ASBE~21 text requires
both practical ability throughout use and economic benefit for a substantive
substitution right. It specifies a presumption when that right is difficult to
determine. Applying that conditional rule requires its stated circumstances;
it does not transform missing lessor information into a reported fact. Nor may
a translator assign arbitrary numerical thresholds to ``substantive'',
``probable'' or ``highly probable''.

\section{Three meanings of disclosure coverage}

The Apple annual/interim preparation inventories 125 quantities, 13 annual notes,
10 interim notes, 70 exhibit references and seven further package obligations.
All 23 notes retain unresolved semantic content; none is recorded as fully
formalised. Downloading their pages answers a useful but narrower question
than constructing a complete company model.

Three populations make that distinction inspectable. Retained-unit coverage
partitions acquired documents into nonoverlapping pages or smaller leaf units,
including duplicates and unreadable material. Assertion-review coverage counts
the separately enumerated material assertions and their review decisions.
Query-exposure coverage identifies the sources, legal facts, methods and
authority that a particular result requires. A fully downloaded annual report
can have incomplete assertion review, and a reviewed debt total can leave a
refinancing query unresolved because an incorporated agreement is missing.
No weighted average of these populations establishes completeness.

Preparation begins with declared annual and interim roots and a policy for
following incorporated references. Each edge is recorded as satisfied,
duplicate or cyclic, outside the declared scope with a reason, inaccessible,
or still unexpanded. A bounded traversal saves its unexplored frontier and
budget so it can resume. An unknown external extent remains unknown. Acquired
proxy and registration documents improve the Apple record, but do not prove
that every executed amendment, issuance term and exhibit dependency has been
closed. The present real package remains incomplete for that reason.

\section{Displayed precision does not determine a rounding rule}

Suppose a statement displays 100 in millions. The scale establishes how the
literal is read; it does not establish how the issuer rounded its underlying
amount. Under an explicitly admitted nearest-rounding assumption, an integer
display $m$ at scale $s$ permits the conservative compatibility enclosure
\[
       s(m-\tfrac12)\leq y\leq s(m+\tfrac12).
\]
The closed endpoints accommodate an unknown tie rule within that nearest-rounding
model. A specified tie rule can exclude an endpoint. Truncation of a positive
amount instead permits $sm\leq y<s(m+1)$. Thus the illustrative display 100
corresponds to 99.5--100.5 million under the first assumption and 100--101 million,
excluding 101, under the second. These are different restrictions on a history.

The retained Apple arithmetic uses the first enclosure as a conditional
compatibility assumption; the issuer's rounding convention has not been
established. Without that assumption or source support, the record states
known display precision and unknown rounding convention and supplies no hard
half-unit constraint. A one-unit difference between a displayed total and the
sum of displayed components is consequently not enough to establish a
contradiction. The source value is scaled once, while its literal and precision
metadata remain available for inspection.

\section{Public availability and system possession are different views}
\label{sec:time-admission}

Northstar's correction explains why a result needs a reporting vintage. A
historical forecast also needs to establish when the exact document bytes were
available. A report downloaded today may carry an old date on its current web
page; that does not prove that the same bytes were public on that date. Public
availability evidence and an ingestion receipt answer different questions.

An investor-information query supplies a cutoff $K$. Each admitted source must
have supported public availability no later than $K$ and evidence linking the
selected bytes to that availability. Later ingestion is allowed. Later
extraction review can read eligible old bytes, but an interpretation that uses
later information fails the cutoff on its own dependencies. This view establishes
public availability, not actual investor awareness or the ability to trade.

A system-as-of query instead supplies an instant $S$ and requires an immutable
receipt showing possession of those bytes by $S$. Publication cannot replace
that receipt. An accepted-layer query also requires review decisions visible
by $S$; a raw-layer query may return a then-unreviewed document. A query asking
both what was public and what the system possessed must pass both filters.
Unknown publication alone need not block a pure possession question.

Archival exploration permits inspection of retained sources with incomplete
timing. With a cutoff, unknown publication yields an unresolved candidate,
not a cutoff-qualified selection. Without a cutoff, the archive may display
the retained lineage, explicitly without a historical-knowledge claim. A known
later version is excluded from an earlier view and cannot suppress an eligible
old assertion. An unresolved potentially relevant version keeps the affected
answer unresolved even when an older assertion is displayed.

Date-only evidence requires particular care. For publication date $D$ in a
named timezone, let $[a_D,b_D)$ be the interval from local midnight to the next
local midnight. A cutoff before $a_D$ excludes the source; a cutoff inside the
interval leaves availability unresolved; a cutoff at or after $b_D$ supports
admission on this date evidence. No arbitrary noon or fixed 24-hour duration
is inserted. Missing timezone leaves the instant comparison unresolved.

\begin{figure}[htbp]
\centering
\begin{tikzpicture}[x=1cm,y=1cm,>=Stealth]
\draw[->] (0,0)--(12,0);
\draw[line width=2pt] (3,0)--(8,0);
\draw (3,-.15)--(3,.15) (8,-.15)--(8,.15);
\node[below,align=center] at (3,-.2) {local midnight\\$a_D$};
\node[below,align=center] at (8,-.2) {next midnight\\$b_D$};
\node[above] at (1.3,.1) {excluded};
\node[above] at (5.5,.1) {time unresolved};
\node[above] at (10.1,.1) {date supports admission};
\end{tikzpicture}
\caption{A dated publication supports historical admission only after its
possible publication interval has ended. The line marks uncertainty about an
instant, not a delay imposed on an already known timestamp.}
\end{figure}

Revision selection follows admission. Original selection returns eligible
roots; latest-available selection returns eligible assertions without an
eligible superseding descendant for the same context. Unrelated roots and
multiple live leaves remain competing assertions. A harmonised view adds
explicit transformations with their own admitted sources, method and opening
bridges. File order never selects a winner.

The synthetic interim example has original first-half revenue 120, revised
first-half revenue 110 and revised annual revenue 260. Compatible subtraction
gives second-half revenue $260-110=150$. The arithmetic $260-120=140$ is easy
but mixes incompatible vintages. An original-view display may show both source
amounts while refusing that derived flow. A change in calendar or perimeter
likewise requires a specific bridge.

The investigated real sources have not yet supplied a qualifying unresolved
same-context contradiction. Known revisions, different reporting bases and
synthetic conflicts remain useful cases of their own kinds. A genuine conflict
case must first match entity, perimeter, concept, dimensions, unit, basis,
period, versions and precision, and investigate extraction and supersession.
The absence of a qualifying example in a bounded search is not evidence that
the corpus contains no conflicts.
